\chapter{Moduli Spaces}\label{moduli-spaces}
\chapterstatus{planned}

Algebraic varieties and schemes that parametrise other geometric objects: representable functors, Hilbert and Quot schemes, quotients by geometric invariant theory, the moduli of curves and of vector
bundles, and the language of stacks. The analytic side (Kodaira--Spencer, Kuranishi) is in Chapter~\ref{deformations}.

\noindent\textbf{Prerequisites.} Chapter~\ref{schemes} (Schemes and Coherent Sheaves), Chapter~\ref{curves} (Algebraic Curves and Riemann Surfaces), Chapter~\ref{algebraic-groups} (Algebraic Groups and
Abelian Varieties) and Chapter~\ref{deformations} (Deformations and Moduli).

\noindent\textbf{Planned contents.}
\begin{enumerate}
\item Moduli problems as functors; fine and coarse moduli spaces; why automorphisms obstruct fine moduli; Grassmannians as the model case.
\item Hilbert and Quot schemes: construction via Castelnuovo--Mumford regularity, tangent spaces and obstructions, Hilbert schemes of points on surfaces.
\item Geometric invariant theory: reductive groups, affine and projective quotients, the Hilbert--Mumford criterion, stability; binary forms and plane cubics as examples.
\item The moduli of curves: $M_g$ via GIT (survey), $\dim M_g = 3g - 3$, stable curves and the Deligne--Mumford compactification $\overline{M}_{g,n}$, $\overline{M}_{0,n}$ in detail.
\item Moduli of vector bundles on curves: stability, Narasimhan--Seshadri (statement), Verlinde formula (survey).
\item Stacks (survey): groupoids, algebraic stacks, $\overline{\mathcal{M}}_{g,n}$ as a Deligne--Mumford stack, intersection theory on it and Witten's conjecture (Kontsevich, statement).
\item Exercises with solutions.
\end{enumerate}
